\section{Attack Thoughts and Cold Hearts}

\begin{quotex}
O you, sweet warm soul, how inspired I was in your arms.
\flright{\textsc{Richard Wagner} in a letter to Judith Gauthier}

\begin{verse}
And when man falls silent in his torment.\\
A god allowed me to say what I'm suffering.
\end{verse}
\flright{\textsc{Goethe}, \emph{Marienbad Elegy}}

\begin{verse}
Und wenn der Mensch in seiner Qual verstummt.\\
Gab mir ein Gott zu sagen, was ich leide.
\end{verse}
\flright{\textsc{Goethe}, \emph{Marienbader Elegie}}

\begin{verse}
I carry peace in me,\\
I carry in myself the forces that strengthen me\\
I want to fulfil myself with these forces of warmth,\\
I want to permeate myself with my will power.\\
And I want to feel how peace pours through all my being,\\
When I am strong, peace as power.\\
To find within me through the power of my striving.
\end{verse}
\flright{\textsc{Rudolf Steiner}}
\end{quotex}


\paragraph{Attack Thoughts}

The recent guest post about correct thinking\footnote{\url{https://gornahoor.net/?p=10120}} coincides synchronistically with our group exercise on observing negative thoughts. On the same day, the quote from the Orthodox priest, the \textbf{Elder Thaddeus of Vitovnica}, about the effect of attack thoughts also showed up.

Although I don't like to promote New Age ideas, the Elder's teaching that our thoughts determine our lives is not so incompatible with the \emph{Course in Miracles}. It is no secret that I diligently performed the exercises daily for one year, more than a couple of decades ago. I did learn the role that our thoughts play in our lives. I certainly don't recommend the main text, but some of the exercises are worthwhile. What is relevant are the exercises on Attack Thoughts\footnote{\url{https://gornahoor.net/?page_id=10130}}.

Conversations and discussions often devolve into mock legal battles. Instead of a dialectic leading toward a higher aim, there is nothing but attack and defence, in an endless loop. If the initial attack is fueled by intense passion, even a guilty plea will often not suffice. Ideally, at least between friends, the attack should not even start. When there is a misunderstanding on one side, or a poorly phrased statement on the other, a request for clarification should be the first step. That most often will rectify the issue before hard divisions are formed.

After all these years, those daily exercises I performed are still fresh in my mind. In a way, this demonstrates the wisdom of Rudolf Steiner's assertion that a mere 5 minutes per day of careful thinking can change your life. Those exercises took little more time than that. I'll leave you with this idea:

\begin{quotex}
Today's idea accurately describes the way anyone who holds attack thoughts in his mind must see the world. Having projected his anger onto the world, he sees vengeance about to strike at him. His own attack is thus perceived as self defence. This becomes an increasingly vicious circle until he is willing to change how he sees. Otherwise, thoughts of attack and counter-attack will preoccupy him and people his entire world. What peace of mind is possible to him then?
\end{quotex}


\paragraph{The Decameron}


Ever since I was a boy, I've been an avid reader. Also, an active reader, like the boy in the \emph{Neverending Story}, who inserts himself into the story. From Tom Corbett to Robert Heinlein novels, then to Ian Fleming and assorted detective and pulp fiction novels. I also read the best authors. But it is in novels like \emph{Siddhartha}, \emph{Steppenwolf}, \emph{Là-bas}, and \emph{Zanoni} that I still find my life story played out.

So the great authors have been my invisible friends who have shaped my life and my thoughts. \textbf{Giovanni Boccaccio} was among them. I regard him as a high initiate of the Fedeli d'Amore, while others only see a wise observer of the human condition, or merely a clever story teller.

His stories are alchemical transformations, that convert the vices and foibles of men and women into a deep insight into humanity, most often by some clever or unexpected catalyst. The stories were intended as a diversion to women since at that time they weren't granted the pastimes of men. He writes regarding women:

\begin{quotex}
Women are restricted by the wishes, whims, or commands of fathers, mothers, brothers, and husbands. They remain most of the time restricted to the narrow confines of their bedrooms, where they sit in apparent idleness, now wishing one thing and now wishing another, turning over in their minds a number of thoughts.
\end{quotex}


But men have more choices:

\begin{quotex}
If men are afflicted by melancholy or ponderous thoughts, they have many ways of alleviating or forgetting them. If they wish, they can take a walk or listen to or look at many different things. They can go hawking, hunting, or fishing. They can ride, gamble, or attend to business.
\end{quotex}


However, in our day, women are free to pursue distractions just as men do. Nonetheless, they still seem to occupy their minds with a number of thoughts. Sometimes, these thoughts take offense at the stories.

Boccaccio anticipated this, and justified his style as an Afterward. Here follow some snippets. Just as Jesus said we are not made impure by what we eat, Boccaccio makes a similar argument:

\begin{quotex}
A corrupt mind never understands a word in a healthy way! And just as fitting words are of no use to a corrupt mind, so a healthy mind cannot be contaminated by words which are not so proper, any more than mud can dirty the rays of the sun.
\end{quotex}


To the women who objected to his words, Boccaccio replied:

\begin{quotex}
Nothing is so indecent that it cannot be said to another person if the proper words are used to convey it; and this, I believe, I have done very well.
\end{quotex}


As for right speech, it all depends on the situation. Boccaccio explains:

\begin{quotex}
These tales were not told in a church, where things must be spoken of with the proper frame of mind and suitable words. Nor were they held in schools of the philosophers, where a sense of propriety is required. But they were told in a garden, in a place suited for pleasure, in the presence of young people, who were nevertheless mature and not easily misled by stories.
\end{quotex}


This blog is my home and my garden and I have invited you all as guests. Sometimes my tales are serious, sometime more humorous or even ribald. Please indulge me as your host. I don't want to offend, but if so, I stand on Boccaccio's defences. If you prefer to convict me instead, things are just how they are.

\paragraph{The Conversion of Abraham}


The reason I include this story is as an antidote to all the hysteria going on about Rome. We can see that the moral level of 14\textsuperscript{th} century Rome was probably worse than what there is today. To not feel dispirited by current events, we need to take the long view of human history. We've seen it all happen before, we've seen empires appear and disappear, battles won and lost, hopes dashed, loves lost, faiths shaken. Yet we persist, knowing the final victory will be ours.

One of the tales was about a Parisian merchant, Giannotto, and his close friend Abraham, who was a Jew. This story is certainly not an attack on Rome. Rather, it shows how the Holy Spirit can transform the worst in man into something Holy. There is no other way to understand it.

Giannotto and Abraham were both respected and trustworthy gentleman. Of course, the topic of religion often came up between these friends of different faiths. Giannotto always failed to convince Abraham to convert.

Eventually Abraham made a proposal. He would travel to Rome to see the Pope and the cardinals. If they were truly holy men, then Abraham would convert. Giannotto was shocked at that since he was aware of the wicked and filthy lives of the clergy in Rome. Believing that Abraham would never convert if he saw that, he tried to dissuade Abraham from his trip, but Abraham persisted.

Once in Rome, Abraham observed carefully the behaviour of the Pope, cardinals, and other prelates. Boccaccio describes what was seen:

\begin{quotex}
They all shamelessly participated in the sin of lust, not only the natural kind of lust but also the sodomitic variety, without the least bit of remorse or shame. And this they did to the extent that the influence of whores and young boys was of no little importance in obtaining great favours. Besides this, he observed that all of them were open gluttons, drinkers, and sots, and that after their lechery, just like animals, they were more servants of their bellies than of anything else; the more closely he observed them, the more he saw that they were all avaricious and greedy for money and that they were just as likely to buy and sell human (even Christian) blood as they were to sell religious objects pertaining to the sacraments or connected to benefices ... they called their blatant simony ``mediation'' and their gluttony ``maintenance'' as if God did not know the intention of these wicked minds.
\end{quotex}


After he returned to Paris, Abraham told Giannotto about his experience:

\begin{quotex}
I don't like them a bit and may God condemn them all; and I tell you this because as far as I was able to determine, I saw there no holiness, no devotion, no good work or exemplary life, or anything else among the clergy. Instead, lust, avarice, gluttony, fraud, envy, pride, and the like and even worse, were so completely in charge there that I believe the city is more a forge for the Devil's work than for God's: in my opinion, the Shepard of yours and all the others as well are trying as quickly as possible and with all the talent and skill they have to reduce the Christian religion to nothing and to drive it from the face of the earth when they really should act as its support and foundation.
\end{quotex}


To Giannotto's surprise and delight, Abraham concluded that, under those circumstance, only the Holy Spirit could be the real foundation. Hence they went off to Notre Dame for Abraham's baptism.


\flrightit{Posted on 2018-09-14 by Cologero}

\begin{center}* * *\end{center}

\begin{footnotesize}
\begin{sffamily}
\texttt{Lavender on 2018-09-15 at 10:01 said:}

I don't understand the ending at all, why is ``only the Holy Spirit could be the real foundation'' the conclusion?

Should we not judge the tree by it's fruit?

\hfill

\texttt{James on 2018-09-15 at 11:30 said:}

I wouldn't confuse the fruits for the fertilizer , Lavender . God's light makes the faithful heart grow . His forgiveness waters its roots . His Spirit breathes life invisibly into each organic heart . I suppose , to learn by analogy , there is a great mystery on as to why such fruitful and sweet things come out of the , literally , crappiest soils .

The clergy are not the fruits of the tree . Can some clergy be fruits ? Yes .

Still , the question is a bit deeper .

What it sounds like to me is that Abraham has had an ``encounter'' with Catholicism already . Perhaps it was through his friend . In essence , Abraham has some idea of something supernatural . He goes to Rome to test if this strange and ``new life'' that he witnesses with his friend or Catholicism as a whole is man-made or God-made .

Seeing how incredibly human the clergy are , he soundly realizes that this unmentioned encounter --- or this intuition that he has that something extraordinary is present (perhaps seeing it through his friend) , he realizes that only a supernatural being could produce such fruits against all ``natural'' odds .

It is an interesting story and idea --- the idea that we can ``eliminate'' the human factor ``through'' the human factor . The miracle of the Church , ironically , is that it literally *is* the body of Christ .

After all , the body of Christ is indeed a human one . His physical body contains all of the imperfections of humanity's physicality . In other words , he was also physically dirty at times . Dirt would cling to Him ; bacteria would help His intestines digest food . It could be pierced . It could bleed . It did get pierced . It did bleed .

The Church is an interesting elevation to a higher consideration of Christ's physicality . The Church becomes a metaphysical meditation on Christ .

And here is where I , at least in my limited knowledge , understand Christianity :

Many Jews could not recognize Christ because they assumed the messiah would be clean and grand .

There is a certain smugness to assume that we are no longer vulnerable to the pitfalls of some of these Jews and yet why do so many of us cling to the idea that for the Church to be true , she must be totally washed clean ?

If the Body of Christ in physicality endured the dirtiness of the Levant's sands , then why do we assume that the spiritual manifestation of this Body in the Church would not suffer the dirt of the world as well ? Yet the very core of the Church never changes even while dirt clings to her . No dogma , teaching , or miracle has been altered or changed even while the lowliest inhabit her halls . After all , despite how dirty the world was in Christ's time , Christ never suffered illness nor was He ever corrupted . Yes , we use the word ``corrupt'' to apply to many clergymen , but the very core --- the interior world of the Church --- her teachings and her spirituality were never infiltrated . It is here that we find ``the tree'' that you are looking for . It is there that you can judge if the fruits exist or not .

I will also say that reading Dante as a young man made it very easy for me to understand the supernatural basis of the Church . The Divine Comedy itself is one such fruit that convinced me of the existence of this Eternal Truth . After all , how could something so ... ``beautiful'' exist ? And yet Dante spends nearly half the lines talking about how ridiculously bad the clergy are . I know Catholics who refuse to read the Sublime poet ``because he puts Popes in Hell'' .

Still , Dante , as Pope Benedict XV would say in an encyclical , is ``the most eloquent singer of the Christian idea'' . Grappling with the paradox and scandal of perfection wrapped in limitation and imperfection (and thus understanding the existential position of the Church in modernity) is the very encounter of the Incarnation . Anyone who is able to eat that flesh and drink that blood (especially as they watched His body being crucified) will truly have life eternal . It is not easy nor is it an intellectual exerise (hence why the flesh must be eaten rather than simply looked upon --- it must be COMMUNED with rather than studied) . Even if one is faithful , only one of the apostles was able to remain at the foot of the cross (and as a reward received Wisdom herself into his home as his mother) .

Anyone who debates this fact rather than peacefully meditates upon it unto acceptance is , in some ways , ``doomed'' . He who lives by the sword dies by it . He who lives by debate (i.e. the sword since the sword is a tool of `bifurcation' or `justice' or `judgement') dies by debate even if he is ultimately right . After all , the apostles who fled and could not synthesize or could not ``reconcile'' GOD DYING ON THE CROSS ended up martyred . This is not some kind of threat for martyrdom is not at all to be looked down upon , but there is a privilege to attaining peaceful acceptance and faith that does not become hindered by internal ``discourse'' . John was able to stay and , thus , John was able to see even to the end of the world . It does not make John ``superiour'' in the earthly sense . He still observed the hierarchy even while taking a ``better'' path . After all , even though the Archangel Michael is the Prince of the Heavenly Host , it is the Seraphim who are ranked higher than even he . (How hierarchy exists in the different emanations is .. not something I know much about anyway so I'm only guessing by intuition here) .

I cannot rightfully express all of these thoughts in just short lines , but there are certainly many threads to this that come to mind .

\hfill

\texttt{argusandphoenix on 2018-09-15 at 12:12 said:}

The Abraham story was helpful: I anticipated the ending. There is something in that.

\hfill

\texttt{Max on 2018-09-16 at 20:16 said:}

There are some relevant thoughts for the discussion in ``The Philosophy of Freedom''. Despite having a difficult time following Steiner's style, I finally managed to finish the book this weekend (along with three other titles that had piled up). His initial remarks together with the concluding ones were the best, while much of the inbetween felt like filling.

``No system of Ethics can ever invent other tasks than the realization of those satisfactions which human desires demand, and the fulfillment of man's moral ideas. No Ethical theory can deprive him of the pleasure which he experiences in the realization of what he desires.''

This is a critique meant for those who would wish for a false start in seemingly safe but bloodless ``oughts'', rather than coming to terms with what is actually found in the soul. Desires are very valuable, and ought not to be squandered. It is not helpful to raise ethical systems that dismisses all pursuit of pleasure as immoral. Without it, humans would ultimately not do much of anything, and even less achieve learning or wisdom. The pleasure associated with fulfilling ones moral ideal, which is the desire of the spirit, makes it worthwhile to endure every pain. Morality means the fullest develpment of human nature. This teaching will however -says Steiner- be easily misunderstood by half-fledged youths and developmentally stunted beings. The mature man makes ues of the generic features and pecularities that he was dealt as a material, and imparts on it a form.

``Anyone who wants to eradicate the pleasure which the fulfillment of human desires brings, will have first to degrade man to the position of a slave who does not act because he wills, but because he must.''

The Philosophy of Freedom sets three distinct options before us; find our motive in our moral imagination, imitate others; do nothing.

Moving on, Steiner cautions us to see too much of the generic woman in women. It risks sounding too much like the cliche that ``everybody is unique'', but the difference for Steiner and other esotericists is that individuality is not just a question of quantification, but essentially belongs to an order of a larger whole. Individuality is not easy to understand as evidenced by its modern distortions.

``Those who always mix their own ideas into their judgement on another person can never attain to the understanding of an individuality.

Women are today still as much victims (to a greater degree than men) to the trap of the generic as at the time of Steiner's reflections, and that despite their supposed emancipation. Instead of ``the general idea of womanhood'', which he laments, we have now progressed to a general idea of whoredom. A definite danger is how desires implanted from the outside comes to be admitted as ones own, obfuscating our true desire, and serving as its parody.

Our contemporaries are also severely mentally challenged in distinguishing between considerations and deeds. It may very well be risky to consider a thought, but neither can we reject it forcefully unless we do. Half-thought thoughts are more dangerous, and hence, do not go there at all, or go all the way.

\hfill

\texttt{Mikkel on 2018-09-17 at 16:58 said:}

I had been researching the Elders as well when Elder Thaddeus specifically came up, perhaps good algorithms at play? Nevertheless I did choose to watch and listen... While many flock to articles that state ``5 quick tips to boost self confidence'', practical techniques and their ideas such as these are largely ignored except in those derivative ways, the gulf feels vast when attempting to discuss the merits of these ideas or even in normal conversation when noticing another's presumptions that their world is built on. Maybe we just need better search results, though an elite will never be formed in that concept... I will journal how they work for me, I hope to make use of it.

On the state of the Church, thank you for sharing that story, interesting to see the variety of responses from the `public': ``close your wallets to the Church until it is reformed'', ``this is what made me stop being Catholic'', it makes me think about the actual type of person who ends at that decision (as well as their allowed influences), as if salvation and the Way can ever be mitigated by external influences. As we seek to understand the Age of the Spirit, it seems relevant that the foundation of the Church is of the Spirit. Joachim of Fiore was mentioned jokingly about being the foundation of a Church that only has one member which he himself created the Church he was in (a recently popular individual was talking about this in an interview), how can ones interior truly reflect the Church of the Spirit though? This is an exercise I've been thinking about recently.

If the type of person like Abraham is out there, perhaps more so there needs to be more Giannotto's to be friends to those Abrahams.

\hfill

\texttt{Mikkel on 2018-09-17 at 17:20 said:}

@Max

It does risk sounding that ``everybody is unique'' (sentimental) and there is reasoning to that. When we look at what Guenon wrote in Man and His Becoming, he specifically denotes the difficulties in the Western perspective interchanging the words Self \& Ego and Personality \& Individuality respectively, specifically critiquing Theosophists' use of the terms interchangeably in the incorrect way, ie. Self/Individuality, Ego/Personality, the reverse is true. Reminds me of why ``Every Man and Every Woman is a star'' is completely misused, when thought of in combination with the next statement ``Every number is infinite; there is no difference.'' this seems to have more significance on the Self and Individuality and the commonality from which it descends from, one of many possibilities but still from the infinite. Quoting Guenon, ``it (center of individuality) effectually possess this reality only through participation in the nature of the `Self', that is, insofar as it is identified therewith by universalitization'', this is from Chapter 2 if curious.

I wonder if seeing ``too much'' of the generic woman in women (or any X in x1 issue) is a issue in projection from that individuality onto the whole in the first place rather than truly understanding the transcendent principle from which it derives from in the first place. Does anyone else remember the Gornahoor posts on mathematical sets and what they mean about the person as analogy? This is what comes to mind reading this. At the same time, that transcendent principle of the Self cannot be individualized as it cannot be fragmented from itself, so perhaps this is the idea that Steiner is speaking to and I think you are saying as well that a one sided approach does very little into really grasping these concepts. To then apply those concepts onto people thinking it is true leads to more ignorance, Matthew 7:1-3?

Of course the emancipation that women have received is nothing of the sort given it has come from a place of revolution, resentment, and the like, to begin with on manufactured `principles'. When you say the dangers of how desires are implanted from the outside, I am reminded of how women and in general people in our society are portrayed in order to be ``valuable'' though they know nothing of what value is, this becomes entrenched in people's brains to the point where nothing else is known any longer and has fully taken over their thought processes, how do we help remove those from ourselves and others?

\hfill

\end{sffamily}
\end{footnotesize}

